\documentclass{article}
\usepackage[utf8]{inputenc}
\usepackage[T1]{fontenc}
\usepackage{graphicx}
\usepackage{algorithm}
\usepackage{algpseudocode}

\usepackage{listings}
\usepackage{xcolor}
\usepackage[margin=2.5cm]{geometry}
\usepackage{float}

\usepackage{amssymb} %Ajoute des symboles mathématiques comme R, Z, forall, exists, etc....

% Les deux lignes suivantes permettent d'enlever le Listing n°_ lorsque l'on fait des programmes
\usepackage{caption}     
\captionsetup[lstlisting]{labelformat=empty}

\lstdefinestyle{monstyle}{
    backgroundcolor=\color{white},
    commentstyle=\color{gray},
    keywordstyle=\color{blue},
    numberstyle=\tiny\color{gray},
    basicstyle=\ttfamily\small,
    breaklines=true,
    numbers=left,
    numbersep=5pt,
    frame=single,
}
\lstset{style=monstyle}


\title{\textbf{Rapport de TD n°2 : Programmation Impérative}}
\author{DERREZ Lorenzo}

\begin{document}
\maketitle

\section{ERRATUM}
\begin{enumerate}
    \item Il y a certains \texttt{main} où il manque des \texttt{return EXIT\_SUCCESS} en fin de fonction.
\end{enumerate}

\section*{\textbf{Exercice 1} : }
\begin{lstlisting}[language=C, caption={Opérations avec \textit{argc} et \textit{argv}}]
    
// Q1 
int main(int argc, char* argv[]){
    int i = 1; 
    while(i<argc){ 
        printf("%s \n", argv[i]);
        i+=1;
    }
    return EXIT_SUCCESS;
}

// Q2

int calc_taille1(int argc, char* argv[]){
    if(argc>1){
        int taille = 0;
        int i = 0; 
        while(argv[1][i] != 0){
            taille +=1;
            i +=1; 
        }
        return taille;
    }
    else{
        return 0;
    }
}


int calc_taille2(int argc, char* argv[]){
    if(argc>1){
        return strlen(argv[1]);
    }
    else{
        return 0;
    }
}

// Q3
int lengthcompare(int argc, char* argv[]){
    if(argc>2){
        int taille1 = strlen(argv[1]);
        int taille2 = strlen(argv[2]);

        if(taille1>taille2){return taille1;}
        else{return taille2;}
    }
    else{
        if(argc==2){
            return strlen(argv[1]);
        }
        else{return 0;}
    }
}

//Q4
int est_voyelle(char caractere){
    if (caractere == 'a' || caractere == 'e' ||
        caractere == 'i' || caractere == 'o' ||
        caractere == 'u' || caractere == 'y' ||
        caractere == 'A' || caractere == 'E' ||
        caractere == 'I' || caractere == 'O' ||
        caractere == 'U' || caractere == 'Y')
    {
        return 1;
    }
    return 0;
}

int nb_voy(char *chaine){
    int n = strlen(chaine);
    int nb = 0;
    int i = 0;

    while (i < n)
    {
        if (est_voyelle(chaine[i]) == 1)
        {
            nb++;
        }
        i++;
    }
    return nb;
}


int main(int argc, char* argv[]){
    int test1 = calc_taille2(argc, argv);
    printf("%d\n", test1);
    return 0;
}

// Test Q3 : 
int main(int argc, char* argv[]){
    int i = 1; 
    while(i<argc){
        int test = nb_voy(argv[i]);
        printf("%d\n", test);
        i+=1; 
    }
    return EXIT_SUCCESS;
}
\end{lstlisting}


\section*{\textbf{Question 2} :}

\begin{lstlisting}[language=C, caption={Implémentationdes fonctions \textit{palindrome, isalpha, tolower}}]
    
// Q1 : 

int est_palindrome(char* chaine){
    int n = strlen(chaine); 
    int i = 0; 
    int parcours = n/2; 

    int cpt = 0; 

    while(i<parcours && (chaine[i] == chaine[n-1-i])){
        i+=1; 
        cpt+=1; 
    }
    
    if(cpt == parcours){
        return 1; 
    }
    return 0; 
}

// Test Q1 

int main(int argc, char** argv){
    int i = 1; 
    while(i < argc){
        int test_i = est_palindrome(argv[i]);
        if (test_i == 1){
            printf("%s est un palindrome.\n", argv[i]);
        }        
        else{
            printf("%s n'est pas un palindrome. \n", argv[i]);
        }
        i+=1; 
    }
    return EXIT_SUCCESS;
}

// Q2 : 
// #include <ctype.h>

// int is_in_alph(char caractere)
// {
//     return isalpha(caractere);
// }

int is_in_alph(char caractere)
{
    if ((caractere >= 'a' && caractere <= 'z') ||
        (caractere >= 'A' && caractere <= 'Z'))
    {
        return 1;
    }
    else
    {
        return 0;
    }
}

// Test Q2 : 
int main(int argc, char** argv){
    int i = 0; 
    while(i<argc){
        int n_i = strlen(argv[i]); 
        int j = 0; 
        while(j<n_i){
            int test_i = is_in_alph(argv[i][j]); 
            if(test_i){
                printf("%c est dans l'alphabet. \n", argv[i][j]);
            }
            else{
                printf("%c n'est pas dans l'alphabet. \n", argv[i][j]);
            }
            j+=1;
        }
        i+=1;
    }
    return EXIT_SUCCESS;
}

// Q3 : 
char maj_to_min(char car)
{
    if (car >= 'A' && car <= 'Z')
    {
        return car + ('a' - 'A');
    }
    else
    {
        return car;
    }
}

int main(int argc, char **argv)
{
    int i = 1;

    while (i < argc)
    {
        int n_i = strlen(argv[i]);
        int j = 0;

        while (j < n_i)
        {
            char test_j = maj_to_min(argv[i][j]);

            printf("%c", test_j);

            j += 1;
        }

        printf("\n");
        i += 1;
    }

    return 0;
}

//Q4 : 

// Amelioration de la fonction palindrome : 
int est_palindrome_update(char* chaine){
    int n = strlen(chaine); 
    int g = 0;
    int d = n-1;

    while(g<d){
        while(is_in_alph(chaine[g]) == 0 && g<d){
            g+=1;
        }
        while(is_in_alph(chaine[d]) == 0 && g<d){
            d-=1;
        }
        if(g>d){
            break;
        }
        if(maj_to_min(chaine[g]) == maj_to_min(chaine[d])){
            d-=1;
            g+=1;
        }
        else{
            return 0;
        }
    }
    return 1;
}


\end{lstlisting}

\section*{\textbf{Exercice 3} :}
\begin{lstlisting}[language=C, caption={Implémentation de la fonction \textit{chaine miroir}}]

//Q1 :

void chaine_miroir(char* chaine){
    int n = strlen(chaine); 
    int i = 0; 
    while(i<n){
        printf("%c", chaine[n-1-i]);
        i+=1;         
    }
    printf("\n");
}

// Autre methode Q1 : 
char* chaine_miroir_update(char* chaine){
    int n = strlen(chaine); 
    int i = 0; 
    char* tab = malloc((n+1)*sizeof(char));
    tab[n] = 0;
    while(i<n){
        printf("%c", chaine[n-1-i]);
        tab[i] = chaine[n-1-i];
        i+=1;         
    }
    printf("\n");
    return tab;
}


int main(int argc, char** argv){
    int i = 1; 
    while(i<argc){
        chaine_miroir(argv[i]); 
        i+=1; 
    }
    return EXIT_SUCCESS;
}

// Q2 : 
// Etant donne que ma fonction de Q1 ne modifie jamais le tableau, je me suis contente de mettre const char* en modification

void chaine_miroir_2(const char* chaine){
    int n = strlen(chaine); 
    int i = 0; 
    while(i<n){
        printf("%c", chaine[n-1-i]);
        i+=1;         
    }
    printf("\n");
}

int main(int argc, char** argv){
    int i = 1; 
    while(i<argc){
        chaine_miroir(argv[i]); 
        i+=1; 
    }
    return EXIT_SUCCESS;
}
//Avec l'autre methode : 

// Autre methode Q1 : 
char* chaine_miroir_2_update(const char* chaine){
    int n = strlen(chaine); 
    int i = 0; 
    char* tab = malloc((n+1)*sizeof(char));
    tab[n] = 0;
    while(i<n){
        printf("%c", chaine[n-1-i]);
        tab[i] = chaine[n-1-i];
        i+=1;         
    }
    printf("\n");
    return tab;
}

\end{lstlisting}

\section*{\textbf{Exercice 4} :}
\begin{lstlisting}[language=C, caption={Implémentation de la fonction \textit{itoa}}]
    
// Q1 : 

void int_to_a(int code){
    int taille = 1; 
    int odg = 1;
    
    while(code>=odg*10){
        odg = odg*10;
        taille+=1;
    }

    char chaine[taille+1]; 
    chaine[taille] = 0;
    taille--; 
    
    int i = 0; 

    while(odg>0){
        char val = '0' + (code/odg);
        chaine[i] = val;
        code = code%odg;
        odg=odg/10;
        printf("%c", val);
        i+=1;
    }
    printf("\n");
}

// Q2 : 

// Extension avec les entiers signes : 
void int_to_a_signed(int code){
    if(code >= 0){
        int_to_a(code);
    }
    else{
        // Il faudrait seulement gerer le cas INT_MIN autrement car cette methode ne fonctionne pas pour lui
        printf("-");
        int_to_a(code-2*code);
    }
}

// Q3 : 
// Test des fonctions : 

// Test Q1 
int main(int argc, char** argv){
    int i = 1; 
    while(i<argc){
        int t_i = atoi(argv[i]);
        int_to_a(t_i);
        i+=1;
    }   
    return EXIT_SUCCESS;
}

// Test Q2 : 
// A COMPLETER // 

\end{lstlisting}

\section*{\textbf{Exercice 5 }:}
\begin{lstlisting}[language=C, caption={Implémentation de la fonction \textit{atoi} pour un tableau d'entiers, puis de caractères}]
    
int* atoi_tab(char* chaine, int *taille){
    int n = strlen(chaine); 
    int i = 0; 

    int cpt = 0;
    int* tab = malloc(n*sizeof(int));
    while(i<n){
        if(chaine[i] != ' '){
            int j = i+1; 
            while(j<n && chaine[j] != ' '){
                j+=1; 
            }
            char chaine_i[j-i+1];
            int l = 0; 
            int temp = i;
            chaine_i[j-i] = 0; 
            while(l<(j-i)){
                chaine_i[l] = chaine[temp];
                l+=1; 
                temp+=1;
            }
            tab[cpt] = atoi(chaine_i);
            cpt++; 
            i = j;
        }
        else{
            i+=1;
        }
    }
    *taille=cpt;
    return tab;
}

void affiche_tab(int* tab, int n){
    int i = 0; 
    while(i<n){
        printf("%d ", tab[i]);
        i+=1;
    }
    printf("\n");
}

int main(int argc, char** argv){
    int i = 1; 
    int taille; 
    while (i<argc){
        int* tab = atoi_tab(argv[i], &taille);
        affiche_tab(tab, taille);
        free(tab);
        i+=1; 
    }
}

// Q2 : 
#include <ctype.h>


int* atoi_tab2(char* chaine, int *taille){
    int n = strlen(chaine); 
    int i = 0; 

    int cpt = 0;
    int* tab = malloc(n*sizeof(int));
    while(i<n){
        if(isdigit(chaine[i])!=0){
            int j = i+1; 
            while(j<n && (isdigit(chaine[j])!=0)){
                j+=1; 
            }
            char chaine_i[j-i+1];
            int l = 0; 
            int temp = i;
            chaine_i[j-i] = 0; 
            while(l<(j-i)){
                chaine_i[l] = chaine[temp];
                l+=1; 
                temp+=1;
            }
            tab[cpt] = atoi(chaine_i);
            cpt++; 
            i = j;
        }
        else{
            i+=1;
        }
    }
    *taille=cpt;
    return tab;
}

int main(int argc, char** argv){
    int i = 1; 
    int taille; 
    while (i<argc){
        int* tab = atoi_tab2(argv[i], &taille);
        affiche_tab(tab, taille);
        free(tab);
        i+=1; 
    }
    Return EXIT_SUCCESS;
}
\end{lstlisting}

\section*{\textbf{Exercice 6} : }

\begin{lstlisting}[language=C, caption={Implémentation du \textit{selection sort} pour des tableaux d'entiers et decaractères}]
    // Q1 : 

//Implementation du tri par selection 

void selection_sort(int* tab, int n){
    int i = 0;

    int indx_ppvl = i;
    int buff_ppl = tab[0];

    while(i<n){
        int j = i+1; 
        indx_ppvl = i;
        buff_ppl = tab[i];
        while(j<n){
            if(tab[j] < buff_ppl){
                buff_ppl = tab[j];
                indx_ppvl = j;
            }
            j+=1; 
        }
        int buff = tab[i]; 
        tab[i] = buff_ppl; 
        tab[indx_ppvl] = buff;
        i+=1;
        
    }
}

// Test Q1 : 

void affiche_tab(int* tab, int n){
    int i = 0; 
    while(i<n){
        printf("%d ", tab[i]);
        i+=1;
    }
    printf("\n");
}

int main(int argc, char** argv){
    int i = 1; 
    int* tab = malloc((argc-1)*sizeof(int));
    while(i<argc){
        int x = atoi(argv[i]);
        tab[i-1] = x; 
        i+=1;
    }
    affiche_tab(tab, argc-1);
    selection_sort(tab, argc-1);
    affiche_tab(tab, argc-1);
    free(tab);
    return EXIT_SUCCESS;
}   

// Q2 : 

void lexico_sort(char** tab, int n){
    int i = 0; 

    int indx_ppl = 0; 
    char* buff_ppl = tab[0];

    while(i<n){
        int j = i+1; 
        indx_ppl = i; 
        buff_ppl = tab[i];
        while(j<n){
            if(strcmp(buff_ppl, tab[j]) > 0){
                buff_ppl = tab[j]; 
                indx_ppl = j; 
            }
            j+=1;
        }
        char* buff = tab[i]; 
        tab[i] = buff_ppl; 
        tab[indx_ppl] = buff; 
        i+=1;
    }
}

void affiche_tab_char(char* chaine){
    int n = strlen(chaine); 
    int i = 0; 
    while(i<n){
        printf("%c",chaine[i]);
        i+=1;
    }
    printf("\n");
}

int main(int argc, char** argv){
    lexico_sort((argv+1), argc-1); // Ici il faut decaler le pointeur pour eviter de prendre en compte ./exo6_exec dans le tri. Il faut alors diminuer la taille de argc. 
    int i = 1; 
    while(i<argc){
        affiche_tab_char(argv[i]); 
        i+=1;
    }
    return EXIT_SUCCESS;
}

\end{lstlisting}

\end{document}